\documentclass[a4paper,11pt]{article}
\usepackage[utf8]{inputenc}
\usepackage[T1]{fontenc}
\usepackage[ngerman]{babel}
\usepackage{amsmath,amssymb}
\usepackage{xcolor}
\usepackage{colortbl}
\usepackage{array}
\usepackage{tikz}
\usepackage{enumitem}
\usepackage{geometry}
\geometry{a4paper,margin=2cm}
\definecolor{bgdark}{HTML}{1E1E1E}
\definecolor{fgtext}{HTML}{E6E6E6}
\definecolor{accent}{HTML}{6CB6FF}
\definecolor{accent2}{HTML}{F2A65A}
\definecolor{gridcol}{HTML}{4A4A4A}
\definecolor{linecol}{HTML}{8A8A8A}
\pagecolor{bgdark}
\color{fgtext}
\arrayrulecolor{linecol}
\newcommand{\karo}[1]{\par\noindent\begin{tikzpicture}\draw[step=5mm, gridcol, very thin] (0,0) grid (\linewidth,#1);\end{tikzpicture}\par}
\newcommand{\aufgabe}[2]{\medskip\par\noindent{\color{accent}\large\textbf{Aufgabe #1}}\hfill{\color{accent2}\small #2}\par\vspace{0.3em}\hrule height0.4pt\vspace{0.6em}}
\newcommand{\kopf}[1]{\noindent{\color{accent}\Large\textbf{Numerik — #1}}\par\vspace{0.4em}\noindent\textbf{Name:}~\rule[-2pt]{5cm}{0.4pt}\hfill\textbf{Datum:}~\rule[-2pt]{3cm}{0.4pt}\par\vspace{0.3em}\hrule height0.8pt\vspace{1em}}
\setlength{\parindent}{0pt}
\begin{document}
\kopf{Fehleranalyse}

\textbf{Notation der Vorlesung:} $f(\cdot)$ bezeichnet die \emph{analytische} (exakte) Lösung des
Problems, $\hat f(\cdot)$ den \emph{numerischen} Algorithmus (Hut). $x$ ist die exakte Eingabe,
$\tilde x$ die gestörte Eingabe (Tilde). Die vier Fehlermaße:
\[
\kappa=|f(x)-f(\tilde x)|,\quad
|\hat f(\tilde x)-\hat f(x)|\le s\,|\tilde x-x|,\quad
|\hat f(x)-f(x)|\le c,\quad
|\hat f(\tilde x)-f(x)|.
\]

\aufgabe{1}{leicht}
Ordnen Sie jede der folgenden verbalen Beschreibungen genau einem der vier Maße
\emph{Kondition $\kappa$}, \emph{Stabilität $s$}, \emph{Konsistenz $c$} oder \emph{Konvergenz} zu:
\begin{enumerate}[label=(\alph*)]
  \item ``Eine kleine Störung der Eingabe verändert die exakte Lösung des Problems nur wenig.''
  \item ``Der Algorithmus liefert bei exakter Eingabe nahezu den exakten Lösungswert.''
  \item ``Eine kleine Störung der Eingabe verändert das vom Algorithmus berechnete Ergebnis nur wenig.''
  \item ``Bei gestörter Eingabe weicht das Algorithmus-Ergebnis insgesamt nur wenig vom exakten Wert ab.''
\end{enumerate}
Erklären Sie außerdem in einem Satz den Unterschied zwischen dem \textbf{Hut} ($\hat f$) und der
\textbf{Tilde} ($\tilde x$).
\karo{4cm}

\aufgabe{2}{leicht--mittel}
Gegeben sei die analytische Funktion $f(x)=x^2$. Untersuchen Sie die Kondition
$\kappa=|f(x)-f(\tilde x)|$ bei der festen absoluten Eingabestörung $\tilde x = x+0{,}1$
an zwei Stellen:
\begin{enumerate}[label=(\alph*)]
  \item an der gut-konditionierten Stelle $x=0{,}5$,
  \item an der schlecht-konditionierten Stelle $x=5$.
\end{enumerate}
Berechnen Sie $\kappa$ jeweils und entscheiden Sie, welche Stelle besser konditioniert ist.
\karo{4cm}

\aufgabe{3}{mittel}
Betrachten Sie die analytische Funktion
\[
f(x)=\frac{1}{1-x}.
\]
Berechnen Sie die Kondition $\kappa=|f(x)-f(\tilde x)|$ für die Störung $\tilde x=x+0{,}01$
an der Stelle $x=0{,}9$ (nahe an der Singularität $x=1$) und zum Vergleich an der Stelle $x=0$.
Quantifizieren Sie, um welchen Faktor das Problem nahe $x=1$ schlechter konditioniert ist.
\karo{5cm}

\aufgabe{4}{mittel}
Die Funktion $f(x)=\sin(x)$ werde durch eine \emph{stückweise konstante} Approximation $\hat f$
mit Schrittweite $h$ angenähert: $\hat f(x)=\sin\!\big(h\cdot\lfloor x/h\rfloor\big)$
(linker Stützpunkt des Intervalls). Betrachten Sie die Eingabe $x=1{,}3$ und die gestörte
Eingabe $\tilde x = x+0{,}25 = 1{,}55$. Bestimmen Sie den Stabilitätsfaktor
\[
s=\frac{|\hat f(\tilde x)-\hat f(x)|}{|\tilde x-x|}
\]
für (a) $h=1$ und (b) $h=0{,}2$. Erklären Sie, warum die kleinere Schrittweite zu einem größeren $s$ führt.
\karo{5cm}

\aufgabe{5}{mittel}
Für $f(x)=\sin(x)$ und dieselbe stückweise konstante Approximation
$\hat f(x)=\sin\!\big(h\cdot\lfloor x/h\rfloor\big)$ aus Aufgabe 4 betrachten Sie die Konsistenz
$c=|\hat f(x)-f(x)|$ an der Stelle $x=1{,}3$.
Berechnen Sie $c$ für (a) $h=0{,}5$ und (b) $h=0{,}25$ und zeigen Sie, dass die kleinere
Schrittweite die Konsistenz verbessert.
\karo{4cm}

\aufgabe{6}{mittel--schwer}
Verwenden Sie wieder $f(x)=\sin(x)$ mit $\hat f(x)=\sin\!\big(h\cdot\lfloor x/h\rfloor\big)$,
$h=0{,}5$, exakte Eingabe $x=1{,}3$ und gestörte Eingabe $\tilde x=1{,}55$.
\begin{enumerate}[label=(\alph*)]
  \item Berechnen Sie den Gesamtfehler (Konvergenzmaß) $|\hat f(\tilde x)-f(x)|$.
  \item Berechnen Sie getrennt die Stabilitätsdifferenz $|\hat f(\tilde x)-\hat f(x)|$ und die
        Konsistenz $|\hat f(x)-f(x)|$ und prüfen Sie die Dreiecksungleichung
        $|\hat f(\tilde x)-f(x)|\le |\hat f(\tilde x)-\hat f(x)|+|\hat f(x)-f(x)|$.
  \item Erläutern Sie den Begriff \emph{Bias}: Wogegen konvergiert $\hat f$ bei festem $h$,
        wenn die Eingabestörung gegen Null geht?
\end{enumerate}
\karo{6cm}

\aufgabe{7}{leicht--mittel}
Beantworten Sie kurz mit Begründung:
\begin{enumerate}[label=(\alph*)]
  \item Warum ist die Kondition $\kappa$ \emph{unabhängig} vom verwendeten Verfahren, die Stabilität $s$
        dagegen \emph{nicht}?
  \item Geben Sie ein Szenario an, in dem ein Problem \emph{gut konditioniert}, das verwendete Verfahren
        aber \emph{instabil} ist. Was folgt daraus für die Brauchbarkeit des Ergebnisses?
\end{enumerate}
\karo{5cm}

\aufgabe{8}{schwer}
Vollständige Fehleranalyse. Gegeben sei $f(x)=\sqrt{x}$, approximiert durch die
\emph{stückweise konstante} Regel $\hat f(x)=\sqrt{\,h\cdot\lfloor x/h\rfloor\,}$ mit $h=0{,}5$.
Exakte Eingabe $x=4{,}2$, Störung $\delta=0{,}4$, also $\tilde x = 4{,}6$.
Berechnen Sie nacheinander mit Zahlenwerten:
\begin{enumerate}[label=(\alph*)]
  \item die Kondition $\kappa=|f(x)-f(\tilde x)|$,
  \item den Stabilitätsfaktor $s=|\hat f(\tilde x)-\hat f(x)|/|\tilde x-x|$,
  \item die Konsistenz $c=|\hat f(x)-f(x)|$,
  \item den Gesamtfehler (Konvergenz) $|\hat f(\tilde x)-f(x)|$,
\end{enumerate}
und beurteilen Sie abschließend, ob Problem und Verfahren an dieser Stelle ``gutartig'' sind.
\karo{6cm}

\end{document}
